%% 
%% Copyright 2007-2020 Elsevier Ltd
%% 
%% This file is part of the 'Elsarticle Bundle'.
%% ---------------------------------------------
%% 
%% It may be distributed under the conditions of the LaTeX Project Public
%% License, either version 1.2 of this license or (at your option) any
%% later version.  The latest version of this license is in
%%    http://www.latex-project.org/lppl.txt
%% and version 1.2 or later is part of all distributions of LaTeX
%% version 1999/12/01 or later.
%% 
%% The list of all files belonging to the 'Elsarticle Bundle' is
%% given in the file `manifest.txt'.
%% 
%% Template article for Elsevier's document class `elsarticle'
%% with harvard style bibliographic references

\documentclass[preprint,12pt]{elsarticle}

%% Use the option review to obtain double line spacing
%% \documentclass[preprint,review,12pt]{elsarticle}

%% Use the options 1p,twocolumn; 3p; 3p,twocolumn; 5p; or 5p,twocolumn
%% for a journal layout:
%% \documentclass[final,1p,times]{elsarticle}
%% \documentclass[final,1p,times,twocolumn]{elsarticle}
%% \documentclass[final,3p,times]{elsarticle}
%% \documentclass[final,3p,times,twocolumn]{elsarticle}
%% \documentclass[final,5p,times]{elsarticle}
%% \documentclass[final,5p,times,twocolumn]{elsarticle}

%% For including figures, graphicx.sty has been loaded in
%% elsarticle.cls. If you prefer to use the old commands
%% please give \usepackage{epsfig}

%% The amssymb package provides various useful mathematical symbols
\usepackage{amssymb}
%% The amsthm package provides extended theorem environments
%% \usepackage{amsthm}

%% The lineno packages adds line numbers. Start line numbering with
%% \begin{linenumbers}, end it with \end{linenumbers}. Or switch it on
%% for the whole article with \linenumbers.
%% \usepackage{lineno}

\usepackage{multirow}
\usepackage{hyperref}
\usepackage{graphics}
\usepackage{booktabs}
\usepackage{multirow}
\usepackage{caption}
\usepackage{amsmath}
\usepackage{siunitx}

\journal{Journal of Hydrology}

\begin{document}

\begin{frontmatter}

%% Title, authors and addresses

%% use the tnoteref command within \title for footnotes;
%% use the tnotetext command for theassociated footnote;
%% use the fnref command within \author or \address for footnotes;
%% use the fntext command for theassociated footnote;
%% use the corref command within \author for corresponding author footnotes;
%% use the cortext command for theassociated footnote;
%% use the ead command for the email address,
%% and the form \ead[url] for the home page:
%% \title{Title\tnoteref{label1}}
%% \tnotetext[label1]{}
%% \author{Name\corref{cor1}\fnref{label2}}
%% \ead{email address}
%% \ead[url]{home page}
%% \fntext[label2]{}
%% \cortext[cor1]{}
%% \affiliation{organization={},
%%             addressline={},
%%             city={},
%%             postcode={},
%%             state={},
%%             country={}}
%% \fntext[label3]{}

\title{Issuing flood warnings at a continental scale: the case of the European Flood Awareness System}

%% use optional labels to link authors explicitly to addresses:
%% \author[label1,label2]{}
%% \affiliation[label1]{organization={},
%%             addressline={},
%%             city={},
%%             postcode={},
%%             state={},
%%             country={}}
%%
%% \affiliation[label2]{organization={},
%%             addressline={},
%%             city={},
%%             postcode={},
%%             state={},
%%             country={}}

\author[inst1]{Jesús Casado-Rodríguez}
\author[inst1]{Peter Salamon}
\author[inst2]{Stefania Grimaldi}
\author[inst2]{Corentin Carton De Wiart}
\author[inst2]{Christel Prudhomme}
\author[inst2]{Calum Baugh}
\author[inst2]{Ervin Zsoter}
\author[inst3]{Ilias Pechlivanidis}
\author[inst3]{Nina Bosshard}
\author[inst4]{Michaela Mikulickova}

\affiliation[inst1]{organization={European Commission, Joint Research Centre, Directorate E - Space, Security and Migration},%Department and Organization
            addressline={Via Enrico Fermi 2749}, 
            city={Ispra},
            postcode={21027}, 
            state={(VA)},
            country={Italy}}

\affiliation[inst2]{organization={European Centre for Medium Range Weather Forecasts},%Department and Organization
            addressline={Shinfield Park}, 
            city={Reading},
            postcode={WC2R 2LS}, 
            country={UK}}

\affiliation[inst3]{organization={Swedish Meteorological and Hydrological Institute},%Department and Organization
            city={Norrköping},
            country={Sweden}}

\affiliation[inst4]{organization={Slovak Hydrometeorological Institute},%Department and Organization
            city={Bratislava},
            country={Slovakia}}
			
\begin{abstract}
%% Text of abstract
The European Flood Awareness System (EFAS) of the Copernicus Emergency Management Service is an operational forecasting system whose aim is to raise awareness about floods in European transnational rivers. It produces probabilistic, medium-range discharge forecasts twice a day by running the open-source hydrological model LISFLOOD with four different meteorological forcing, two deterministic forecasts from the DWD (German Weather Service) and the ECMWF (European Centre for Medium Range Weather Forecasts), respectively, and two probabilistic forecasts from ECMWF and the Cosmo Consortium (COSMO-LEPS). Based on these forecasts, flood notifications are issued to the EFAS partners if a set of criteria is met: contributing area larger than 2000 km², lead time from 2 to 10 days, at least one deterministic model exceeds the discharge threshold (5-year return period), and at least one probabilistic model predicts 30\% exceedance probability of that discharge threshold for three or more consecutive forecasts. 
However, the operational EFAS is being regularly updated, so the configuration of EFAS has changed since the time these notification criteria were defined. For instance, the temporal resolution has increased from daily to 6-hourly, and the spatial resolution is planned to improve from 5km to approximately 1.5 km (1 arc-minute). 
This study aims at assessing the skill of the notification criteria above presented with the current system setup, and to derive a new set of criteria that optimizes the notification skill. We will focus on three research questions: (i) how can we combine the different models (deterministic and probabilistic) into a grand ensemble and what probability threshold optimizes skill? (ii) Is the persistence criterion (i.e. 3 consecutive forecasts need to provide persistent predictions of high flood risk) adding to the skill both at shorter and larger lead times? (iii) Can we reduce the contributing area threshold without compromising skill?
The study will make use of reanalysis, driven by meteorological observations, and forecast data at over 1200 points across Europe. By comparing the reanalysis data with the simulated discharge threshold, a total of 678 “observed” flood events have been identified in the years 2021 and 2022. We have tested 4 approaches in which to combine the 4 meteorological forcing and several combinations of the probability and persistence criteria. For each of the runs we computed binary classification skill metrics such as recall, precision and the f-beta score, and we have searched for the combination of criteria that maximizes this latter metric. We have also looked into how the skill of the system evolves with lead-time and catchment area.
The outcome of this study will be applied to the EFAS operational system, directly impacting the preparedness of the relevant authorities in future flood events.
\end{abstract}

%%Graphical abstract
%\begin{graphicalabstract}
%\includegraphics{grabs}
%\end{graphicalabstract}

%%Research highlights
%\begin{highlights}
%\item Research highlight 1
%\item Research highlight 2
%\end{highlights}

\begin{keyword}
%% keywords here, in the form: keyword \sep keyword
EFAS \sep flood warnings \sep floods early warning system \sep probabilistic forecasting
%% PACS codes here, in the form: \PACS code \sep code
\PACS 0000 \sep 1111
%% MSC codes here, in the form: \MSC code \sep code
%% or \MSC[2008] code \sep code (2000 is the default)
\MSC 0000 \sep 1111
\end{keyword}

\end{frontmatter}

%% \linenumbers

%% main text
\section{Sample Section Title}
\label{sec:sample1}

%% For citations use: 
%%       \citet{<label>} ==> Jones et al. [21]
%%       \citep{<label>} ==> [21]
%%

Lorem ipsum dolor sit amet, consectetur adipiscing \citep{Fabioetal2013} elit, sed do eiusmod tempor incididunt ut labore et dolore magna \citet{Blondeletal2008} aliqua. Ut enim ad minim veniam, quis nostrud exercitation ullamco laboris nisi ut aliquip ex ea commodo consequat. Duis aute irure dolor in reprehenderit in voluptate velit esse cillum dolore eu fugiat nulla pariatur. Excepteur sint occaecat cupidatat non proident, sunt in culpa qui officia deserunt mollit \citep{Blondeletal2008,FabricioLiang2013} anim id est laborum.

Lorem ipsum dolor sit amet, consectetur adipiscing elit, sed do eiusmod tempor incididunt ut labore et dolore magna aliqua. Ut enim ad minim veniam, quis nostrud exercitation ullamco laboris nisi ut aliquip ex ea commodo consequat. Duis aute irure dolor in reprehenderit in voluptate velit esse cillum dolore eu fugiat nulla pariatur. Excepteur sint occaecat cupidatat non proident, sunt in culpa qui officia deserunt mollit anim id est laborum see appendix~\ref{sec:sample:appendix}.

%% The Appendices part is started with the command \appendix;
%% appendix sections are then done as normal sections
\appendix

\section{Sample Appendix Section}
\label{sec:sample:appendix}
Lorem ipsum dolor sit amet, consectetur adipiscing elit, sed do eiusmod tempor section \ref{sec:sample1} incididunt ut labore et dolore magna aliqua. Ut enim ad minim veniam, quis nostrud exercitation ullamco laboris nisi ut aliquip ex ea commodo consequat. Duis aute irure dolor in reprehenderit in voluptate velit esse cillum dolore eu fugiat nulla pariatur. Excepteur sint occaecat cupidatat non proident, sunt in culpa qui officia deserunt mollit anim id est laborum.

%% If you have bibdatabase file and want bibtex to generate the
%% bibitems, please use
%%
\bibliographystyle{elsarticle-num-names} 
\bibliography{cas-refs}

%% else use the following coding to input the bibitems directly in the
%% TeX file.

% \begin{thebibliography}{00}

% %% \bibitem[Author(year)]{label}
% %% Text of bibliographic item

% \bibitem[ ()]{}

% \end{thebibliography}
\end{document}

\endinput
%%
%% End of file `elsarticle-template-num-names.tex'.
